% ID: SEL-D0065-OA-E01
% ESTADO: validada. S065.
\textbf{a)} El recinto está sobre $y=0$ y debajo de $y=3-x$ y $y=3+x$, con $x\le2$. Las fronteras se incluyen.
\begin{center}
\setlength{\unitlength}{0.9mm}
\begin{picture}(95,54)(-48,-9)
\qbezier(-43,0)(-2,0)(39,0)
\qbezier(0,0)(0,20)(0,40)
\qbezier(-36,0)(-18,18)(0,36)
\qbezier(0,36)(12,24)(24,12)
\qbezier(24,12)(24,6)(24,0)
\put(40,0){$x$}\put(2,41){$y$}
\put(-42,-6){\scriptsize $(-3,0)$}\put(-13,38){\scriptsize $(0,3)$}
\put(25,12){\scriptsize $(2,1)$}\put(22,-6){\scriptsize $(2,0)$}
\end{picture}
\end{center}
\textbf{b)} Las rectas oblicuas se cortan en $(0,3)$. El extremo izquierdo del eje horizontal es $(-3,0)$. Sobre $x=2$, la región va de $(2,0)$ a $(2,1)$.

\textbf{c)} En $(-3,0),(0,3),(2,1),(2,0)$, los valores de $F$ son $6,-3,-5,-4$. Por linealidad, el máximo es $\boxed{6}$ sólo en $(-3,0)$ y el mínimo es $\boxed{-5}$ sólo en $(2,1)$.
